\chapter{Introducción}
\label{cap:introduccion}

\section{Contexto}

La pesca en Baja California tiene gran importancia económica y cultural. Con más de 1{,}500
kilómetros de litorales y una biodiversidad marina privilegiada, el estado se ubica entre los
principales productores pesqueros de México tanto en volumen como en valor, y sus flotas se
componen principalmente de embarcaciones menores de pesca ribereña. En 2016, Baja California
registró la producción de 45 especies entre pesca y acuacultura, donde sobresalen la sardina,
el atún, el erizo, la langosta, el calamar y el camarón. Esta actividad económica genera
numerosos empleos en las cadenas de captura, procesamiento y distribución. No obstante, el
sector enfrenta problemas por la incertidumbre en el otorgamiento de permisos, la pesca ilegal
y la insuficiente infraestructura.

La organización no gubernamental COBI (Comunidad y Biodiversidad, A.C.) colabora con las
cooperativas de la región y reconoce la importancia de los datos oceanográficos ---temperatura
del mar, pH, salinidad, viento, marea--- para comprender la dinámica de las especies que se
capturan. Anticipar el volumen de captura de la siguiente temporada le permite a una
cooperativa dimensionar insumos, tripulación y compromisos de venta: es una decisión operativa
concreta, con un horizonte natural de unos tres meses, que hoy se toma con la experiencia
acumulada de los pescadores y sin ninguna herramienta cuantitativa.

Este trabajo construye esa herramienta. Se concentra en la langosta roja
(\emph{Panulirus interruptus}) porque es la especie con el registro de arribos más largo y
consistente en la zona, y la extiende a abulón y erizo cuando los datos lo permiten.

\section{Planteamiento del problema}

El problema tiene tres dificultades que lo hacen algo más que un ejercicio estándar de
pronóstico de series de tiempo.

\begin{enumerate}
  \item \textbf{Escasez de datos.} La unidad de análisis útil no es ``la pesca de Baja
  California'' sino la serie \emph{especie $\times$ \ue{}}: cada cooperativa tiene su propia
  zona de captura y su propio comportamiento. Al desagregar así, cada serie se queda con del
  orden de $10^{2}$ a $10^{3}$ observaciones útiles ---unas pocas temporadas--- que es uno o
  dos órdenes de magnitud menos de lo que los modelos modernos de pronóstico suelen requerir.

  \item \textbf{Cambios de régimen ambiental.} La serie de San Quintín registra una caída
  abrupta a partir de la temporada 2021-2022 que ningún modelo del planteamiento inicial
  anticipó, y que todos sobre-predijeron. Un modelo entrenado sobre el régimen previo
  extrapola un nivel de captura que ya no existe.

  \item \textbf{Un pronóstico puntual no es lo que la decisión necesita.} La captura se
  concentra en pocos días de apertura de temporada, de modo que el error diario de cualquier
  modelo es grande aun cuando el agregado de la temporada sea razonable. Para decidir cuántos
  insumos comprar, un rango con cobertura conocida es más útil ---y más honesto--- que un
  número puntual acompañado de una métrica de error promedio.
\end{enumerate}

\section{Pregunta de investigación e hipótesis}

La pregunta que guía el trabajo es: \textbf{¿cómo construir un pronóstico del volumen de
captura por temporada que sea útil para una cooperativa pesquera, dado el volumen de datos
realmente disponible?} De ella se desprenden tres preguntas subordinadas, que ordenan el
documento:

\begin{enumerate}
  \item ¿Qué mecanismo explica los años atípicos, y puede incorporarse como variable
  explicativa?
  \item ¿Es el factor limitante la sofisticación del modelo o el volumen de datos?
  \item ¿Puede el pronóstico entregarse con una medida de incertidumbre calibrada y verificable?
\end{enumerate}

La \textbf{hipótesis central} es que el factor limitante es el \emph{volumen de datos} ---el
número de temporadas y de series disponibles--- y no la capacidad del modelo. La hipótesis es
falsable, y el trabajo la somete a prueba en ambas direcciones: aumentando los datos (de 2 a
21 series de langosta y de $\sim$3 a $\sim$9 temporadas) y aumentando la capacidad del modelo
(hasta un Transformer de fusión temporal), para observar cuál de las dos palancas mueve el
resultado.

\section{Objetivos}

\noindent\textbf{Objetivo general.} Construir y evaluar un sistema de pronóstico del volumen
diario de captura por especie y \ue{}, con incertidumbre calibrada, reproducible y entregable
a las cooperativas pesqueras de Baja California.

\noindent\textbf{Objetivos específicos.}
\begin{enumerate}
  \item Reconstruir de forma reproducible y probada el proceso de datos, uniendo las fuentes
  de arribos disponibles (COBI y CONAPESCA) con covariables oceanográficas satelitales.
  \item Implementar un índice de olas de calor marinas según \citet{hobday2016} y evaluar si
  explica los años atípicos de captura.
  \item Comparar modelos estadísticos, de aprendizaje de máquina y de aprendizaje profundo
  sobre una partición temporal estricta, con métricas explícitas.
  \item Evaluar si un modelo global multi-especie y multi-\ue{} transfiere señal entre series.
  \item Producir un pronóstico probabilístico con cobertura empírica verificable.
  \item Contrastar la hipótesis de datos contra una arquitectura de vanguardia (TFT).
  \item Entregar el resultado como un producto operativo consultable.
\end{enumerate}

\section{Aportaciones}

\begin{itemize}
  \item Un \textbf{diagnóstico causalmente informado} del año atípico: el índice de olas de
  calor marinas conecta la caída de 2021-2022 con el régimen cálido de 2019-2021, mecanismo
  documentado para Baja California por \citet{villasenor2024}.
  \item Un \textbf{modelo global} sobre 44 series con la variable objetivo en escala
  logarítmica, que supera o empata a los modelos específicos en cuatro de cinco series.
  \item Un \textbf{pronóstico probabilístico calibrado} por regresión cuantílica
  conformalizada, con cobertura estable incluso durante olas de calor marinas.
  \item Una \textbf{prueba explícita de la hipótesis de datos}, en las dos direcciones:
  añadiendo series y temporadas (sí mejora, y vuelve reproducible la calibración) y añadiendo
  capacidad de modelo (no mejora).
  \item Un \textbf{producto operativo} desplegable que sirve el intervalo calibrado por
  especie y \ue{}, y un repositorio reproducible con pruebas automatizadas.
\end{itemize}

\section{Estructura del documento}

El Capítulo~\ref{cap:revision} revisa la literatura de pronóstico pesquero, de olas de calor
marinas y de predicción conformal. El Capítulo~\ref{cap:metodologia} plantea formalmente el
problema, la partición temporal, las métricas y los modelos, incluido el procedimiento
conformal. El Capítulo~\ref{cap:datos} describe las fuentes, el proceso de datos y el conjunto
consolidado. El Capítulo~\ref{cap:resultados} presenta los resultados, desde la comparación
inicial de modelos puntuales hasta la prueba de techo con el TFT. El
Capítulo~\ref{cap:producto} describe el producto operativo. El Capítulo~\ref{cap:conclusiones}
concluye, enumera las limitaciones y propone trabajo futuro.
